% ============================================================
%  Preface.tex — Lời nói đầu
% ============================================================
\section*{LỜI NÓI ĐẦU}
\phantomsection
\addcontentsline{toc}{section}{\numberline{}LỜI NÓI ĐẦU}
\thispagestyle{empty}

Sự bùng nổ của \textit{Internet of Things} (IoT) trong những năm gần đây đã tạo
ra một hạ tầng kết nối khổng lồ với hàng tỷ thiết bị nhúng hoạt động trong các
môi trường đa dạng từ nhà thông minh, y tế, đến hệ thống công nghiệp và thành
phố thông minh. Nền tảng định tuyến cho các mạng IoT hạn chế tài nguyên
--- giao thức \textbf{RPL} (RFC~6550) --- đóng vai trò không thể thiếu trong
việc duy trì liên lạc giữa các nút cảm biến. Tuy nhiên, thiết kế nhẹ nhàng
của RPL cũng mang lại nhiều lỗ hổng bảo mật nghiêm trọng mà các biện pháp mặc
định của giao thức chưa đủ sức bảo vệ.

Đồ án môn học này được thực hiện với mục tiêu tìm hiểu, phân tích bốn dạng
tấn công tiêu biểu nhằm vào RPL --- DIS Flooding, Blackhole, Decreased Rank và
VeRA/DODAG Version Number Attack --- đồng thời nghiên cứu và đề xuất kiến trúc
\textbf{Hệ thống Phát hiện Xâm nhập} (IDS) có khả năng phát hiện các mối đe
dọa này trong môi trường mô phỏng Contiki-NG/Cooja.

Trong quá trình thực hiện đồ án, chúng tôi xin gửi lời cảm ơn chân thành đến
thầy \textbf{Trần Quang Đức}, người đã tận tình hướng dẫn, góp ý và động viên chúng tôi
hoàn thành công trình này. Chúng tôi cũng xin cảm ơn các thầy cô trong \textbf{Trường
Công nghệ Thông tin và Truyền thông} đã trang bị cho chúng tôi nền tảng kiến thức vững chắc trong suốt
quá trình học tập.

Mặc dù đã cố gắng hết sức, đồ án chắc chắn còn nhiều thiếu sót. Chúng tôi rất mong
nhận được sự góp ý của thầy cô và các bạn để công trình được hoàn thiện hơn.

\vspace{1cm}
\begin{flushright}
Hà Nội, 04/05/2026.

\vspace{0.5cm}
\textbf{Nhóm sinh viên thực hiện}

\vspace{2cm}
\textbf{Trần Hoàng Long} \\
\textbf{Hoàng Công Phát} \\
\textbf{Trần Mạnh Dũng}
\end{flushright}

\cleardoublepage